% !TeX program = pdflatex
%-------------------------
% Cover letter
% Author : Javier Tarazona
% Adapted for: Datadog - Applied Science Intern (Paris)
%------------------------
\documentclass[11pt,a4paper]{article}

\usepackage[utf8]{inputenc}
\usepackage[T1]{fontenc}
\usepackage{lmodern}
\usepackage{geometry}
\usepackage[hidelinks]{hyperref}
\usepackage{xcolor}

\geometry{
  top=1.4cm,
  bottom=1.4cm,
  left=1.8cm,
  right=1.8cm
}

\hypersetup{
  pdftitle={Javier Andres TARAZONA JIMENEZ - Cover Letter - Datadog Applied Science Intern},
  pdfauthor={Javier Andres TARAZONA JIMENEZ},
  pdfsubject={Application for the Applied Science Intern position (Requisition ID 1075)},
  pdfkeywords={Applied Science, Machine Learning, Generative AI, LLM, Computer Vision, Anomaly Detection, Observability}
}

\pagestyle{empty}
\frenchspacing
\setlength{\parindent}{0pt}
\setlength{\parskip}{0.45em}

\begin{document}

\begin{minipage}[t]{0.64\textwidth}
\textbf{\large Javier Andres TARAZONA JIMENEZ}\\
{\small Engineering Student in Artificial Intelligence -- ENSTA, IP Paris}\\
{\small Massy, Île-de-France, France}\\
{\small \href{mailto:javitar06@gmail.com}{javitar06@gmail.com} \textbullet{} +33 7 46 36 97 75}\\
{\small \href{https://www.linkedin.com/in/javier-andres-tarazona-jimenez-84b489222/}{LinkedIn} \textbullet{} \href{https://github.com/JavierTarazona06}{GitHub}}
\end{minipage}
\hfill
\begin{minipage}[t]{0.34\textwidth}
\raggedleft
\textbf{Datadog}\\
{\small Applied Science Hiring Team}\\
{\small Paris, France}
\end{minipage}

\vspace{0.4cm}

{\raggedleft Massy, 30 September 2026\par}

\vspace{0.2cm}

\textbf{Subject: Application for the Applied Science Intern position - Paris}

\vspace{0.1cm}

Dear Hiring Team,

I am an engineering student in a double-degree programme between Universidad Nacional de Colombia and ENSTA, a member school of Institut Polytechnique de Paris, and I am applying for the Applied Science Intern position as my six-month end-of-studies internship, starting in April 2027.

I came to computer science out of curiosity for systems whose apparent simplicity hides considerable complexity, such as the Internet, financial infrastructures or cameras, and out of a wish to turn abstract ideas into systems I could understand and build myself. Artificial intelligence then offered a different paradigm: instead of programming every behaviour explicitly, building systems that extract information from data, model part of reality and support better decisions. Its transversal reach drew me towards applied AI and to ENSTA's AI-focused engineering curriculum in France. My work has since followed three complementary axes: data analysis, in my research internship at ENSTA Paris -- U2IS and in decision-oriented work for the ``Tameo'' autonomous boat; perception from images and text, through Tameo's detection and motion-decision pipeline, lane-marking analysis at Engin A.I and an NLP application built at APEDYS91 to support children with dyslexia; and software development, since AI only creates value once deployed as a working service, as in SharpSight, ORIUN and my work at Engin A.I.

Datadog interests me because its core challenge is a data problem: observing the systems companies depend on and turning a continuous flow of metrics, logs and events into signals engineers can act on. Automatic detection, prediction of incidents and investigation of their causes are research problems with an immediate product consequence. Two examples I studied illustrate this: Faulty Deployment Detection, which analyses new deployments to identify faulty ones and became a production feature, and Watchdog, which applies anomaly detection to streaming data at scale. This role sits at that junction, combining model benchmarking, large-scale data exploration and the transfer of research into customer-facing features, which is closer to how I work than training a model on a fixed dataset.

That research-to-production loop is what I can bring. At U2IS, I did not stop at reading papers: I built a Python/Blender pipeline generating multi-view training data, then pre-trained Vision Transformers in PyTorch for novel view synthesis through curriculum learning, treating the model as one component of a larger pipeline. On Tameo, a team of five I now lead, I benchmarked and fine-tuned detection models on specialised datasets, raising mAP, F1 and precision/recall to 80\% - 90\%, the benchmarking reasoning the role requires. For generative AI, my APEDYS91 application combines rules and local LLMs (Mistral, Gemma) with spaCy-based entity and number checks and new-token-ratio control against hallucinated content, plus model selection adapted to available RAM/VRAM: constraints that matter even more for customer-facing features. Observability data would be new to me; I would approach it with my coursework in probability, statistics and stochastic modelling, and the habit of learning unfamiliar domains quickly.

In return, Datadog would expose me to problems academic projects cannot reproduce: detecting anomalies and explaining their root causes across metrics, events, code and service topology, under the scale and reliability constraints of a production platform. Working alongside experienced engineers and applied scientists, with the support of mentorship, would help me turn a profile rooted in computer vision, with up-to-date but mostly academic machine learning knowledge, into solid practice as an applied scientist or ML engineer in production, a path I would like to pursue with the company beyond the internship. I would welcome the opportunity to discuss how I could contribute to your team.

Yours sincerely,

\vspace{0.2cm}

Javier Andres TARAZONA JIMENEZ

\end{document}
